% Source: https://artofproblemsolving.com/community/c3h11891_toppers_random_nats_problems
% Stable identity: aops-topic-11891-q01
Clearing denominators, $15(m+n)=2mn$, so $4mn-30m-30n=0$ and
\[
(2m-15)(2n-15)=225 .
\]
Since $\dfrac{1}{n}<\dfrac{2}{15}$ we have $n>7$, and similarly $m>7$, so both factors are odd numbers greater than $-1$; as their product is positive, both are positive. With $m>n$, the factor $2n-15$ is a divisor of $225$ below $15$:
\[
\begin{array}{c|c|c|c}
2n-15 & 2m-15 & (n,m) & m \text{ a multiple of } n?\\ \hline
1 & 225 & (8,120) & \text{yes}\\
3 & 75 & (9,45) & \text{yes}\\
5 & 45 & (10,30) & \text{yes}\\
9 & 25 & (12,20) & \text{no}
\end{array}
\]
So $(m,n)=(20,12)$ and $m+n=\boxed{32}$.
